%% ============================================================================
\section{Literature Review}
\label{ch:literature}
%% ============================================================================

This chapter reviews three bodies of work that the thesis sits between:
semantic segmentation architectures (Section~\ref{sec:lit-seg}),
deep learning applied to archaeological prospection (Section~\ref{sec:lit-arch}),
and input-pipeline performance for large-scale training (Section~\ref{sec:lit-io}).
The first two establish the workload; the third establishes the gap that
\pads{} addresses.

\subsection{Semantic segmentation architectures}
\label{sec:lit-seg}

Semantic segmentation assigns a class label to every pixel of an image.
The modern formulation begins with fully convolutional encoder--decoder networks,
of which U-Net~\cite{Ronneberger15} is the canonical instance: a contracting path
captures context, an expanding path recovers spatial resolution, and skip
connections carry high-frequency detail across the bottleneck.
U-Net was introduced for biomedical imaging under severe data scarcity, which is
one reason it transferred well to remote sensing, where annotated corpora are
similarly small.

Two families extend this design in directions relevant here.
The first enlarges the effective receptive field without sacrificing resolution.
DeepLabV3+~\cite{Chen18} combines atrous (dilated) convolution and atrous spatial
pyramid pooling with a decoder module that recovers object boundaries, capturing
multi-scale context at high capacity.
The second inserts attention into the encoder--decoder.
MA-Net~\cite{Fan20} augments a U-Net-like structure with self-attention in the
latent space, allowing the network to weight latent features by content - an idea
imported from the transformer literature~\cite{Vaswani17}.
Developed for liver and tumour segmentation, MA-Net has since been applied to
remote sensing, and it is the architecture selected as best-performing by
Casini et al.~\cite{Casini23}, which is why it is the reproduction baseline here.

A third family replaces the convolutional encoder entirely.
SegFormer~\cite{Xie21} pairs a hierarchical, positional-encoding-free transformer
encoder with a lightweight all-MLP decoder, achieving competitive accuracy at
markedly lower per-step compute than heavy CNN decoders.
For the purposes of this thesis that low compute cost is not incidental: it makes
SegFormer the configuration in which data-loading stalls are most exposed, since a
model that computes quickly spends proportionally more of its time waiting.

Encoder choice is largely orthogonal to decoder architecture.
EfficientNet~\cite{Tan19} obtains strong accuracy-per-parameter through compound
scaling of depth, width and resolution, and EfficientNet-B3 is used throughout this
work as the encoder for both CNN decoders, following the reference study.
All architectures are instantiated through
\code{segmentation\_models\_pytorch}~\cite{Iakubovskii19}, which allows encoder and
decoder to be selected independently and supplies the loss functions used here.

\subsection{Deep learning for archaeological prospection}
\label{sec:lit-arch}

Archaeology is what Bickler terms a ``slow data'' field: annotated examples accrue
through fieldwork, not through scraping, and corpora are correspondingly small.
The standard response is transfer learning - initialising from a model pretrained
on a large general corpus such as ImageNet~\cite{Deng09} and fine-tuning on the
specific task.

Applications divide by output granularity.
Tile classification labels a map tile as containing a site or not; object detection
predicts bounding boxes; semantic segmentation predicts a shape.
Orengo et al.~\cite{Orengo20} applied machine-learning classification to multisensor
and multitemporal satellite data for mound detection, working at the level of
per-pixel classification over large regions.
Trier et al.~\cite{Trier19} used deep networks on airborne laser scanning data for
semi-automatic mapping of archaeological topography, and
Guyot et al.~\cite{Guyot21} combined detection and segmentation of archaeological
structures from LiDAR.
Verschoof-van der Vaart and Lambers~\cite{Verschoof22} report on applying automated
object detection within archaeological practice, emphasising the integration
question rather than the accuracy question alone.
A closely related strand uses object detection rather than segmentation to
localise individual features: Caspari and Crespo~\cite{Caspari19} train a CNN to
detect early Iron Age burial mounds in the Eurasian steppe, and Chen et
al.~\cite{Chen21} adapt off-the-shelf object detectors to find stone kurgans in the
Altai Mountains; both report that the model narrows the area a human surveyor must
inspect rather than replacing survey outright.
Guyot et al.\ illustrate the same architectural shift within a single research
line: an earlier Random Forest classifier over multi-scale LiDAR-derived
topographic features~\cite{Guyot18} is followed, three years later, by a
deep-learning detection-and-segmentation pipeline on the same class of
data~\cite{Guyot21}.

The Mesopotamian floodplain dataset used in this thesis originates from a specific
research programme at the University of Bologna, and that programme's own
publication history illustrates how this literature accumulates.
Roccetti et al.~\cite{Roccetti20} first assessed the feasibility of a deep-learning
model for discovering new sites in this floodplain; Casini et
al.~\cite{Casini22} followed with the preliminary semantic-segmentation study that
this thesis reproduces (Appendix~\ref{app:reproduction}); Casini et
al.~\cite{Casini23} then added the human-in-the-loop adjudication step discussed
below; and, most recently, Pistola et al.~\cite{Pistola25} retrained the same
underlying model on CORONA satellite imagery from the 1960s for the Abu Ghraib
district west of Baghdad, reporting a segmentation IoU above 85\% and confirming
four previously unrecorded sites through field verification.
Two other groups apply CNNs to CORONA imagery in different regions: Soroush et
al.~\cite{Soroush20} detect qanat shafts in the Kurdistan Region of Iraq, and
Vadineanu et al.~\cite{Vadineanu24} use explainability techniques on
CORONA-trained classifiers to delineate mounded settlements in Upper Mesopotamia.
Both confirm that CORONA-based detection is an established, independently
replicated task rather than a peculiarity of this dataset, though neither paper
reports the per-file storage cost of the imagery, which is this thesis's specific
concern with it.

Two themes from this literature bear directly on the present work.
First, evaluation is unusually fraught.
Fiorucci et al.~\cite{Fiorucci22} argue that standard detection measures are poorly
matched to archaeological objects on LiDAR and propose alternatives; Casini et
al.~\cite{Casini23} make the related observation that IoU and site-level detection
accuracy can move in opposite directions, because a more conservative model
produces fewer positive predictions, shrinking the union term and raising IoU while
missing sites and lowering recall.
Any model comparison that relies on IoU alone must be read with that caveat.

Second, labels are noisy in a structured way.
A substantial fraction of catalogued sites are no longer visible on modern imagery,
having been levelled for agriculture or obscured by development; a model that
correctly produces no contour for such a tile is penalised as a false negative
against the annotation.
Casini et al.~\cite{Casini23} address this with human-in-the-loop adjudication,
which raises reported detection accuracy from roughly 63\% to 80\%.
This is a property of the corpus, not of any particular model, and it bounds what
any reproduction can be expected to achieve.

\subsection{Input pipelines and I/O for large-scale training}
\label{sec:lit-io}

The literature above is concerned with what the model computes.
A separate and smaller literature is concerned with whether the GPU is
given anything to compute.
General surveys of GPU-oriented deep-learning optimisation~\cite{Mittal19} cover
compute- and memory-level techniques - kernel fusion, mixed precision,
distributed parallelism - but treat the input pipeline as out of scope; the
work reviewed below fills that gap specifically.

Mohan et al.~\cite{Mohan21} characterise \emph{data stalls} in DNN training and show
that they are common rather than exceptional: across a range of workloads,
accelerators spend a substantial fraction of wall-clock time blocked on fetching and
preprocessing, and the effect worsens as accelerators get faster relative to storage
and host CPUs.
Their system, CoorDL, addresses this through coordinated caching and by eliminating
redundant preprocessing across concurrent jobs.
Murray et al.~\cite{Murray21} describe \code{tf.data} and make the same argument from
the framework side, arguing that input pipelines deserve to be treated as
first-class, optimisable programs rather than as glue code.
Svogor et al.~\cite{Svogor22} profile the PyTorch \code{DataLoader} specifically
against high-latency, S3-like object storage and report that a modified
\code{ConcurrentDataloader} recovers up to a $12\times$ reduction in batch-loading
time - independent confirmation, on different storage and a different framework
layer, that the many-small-file problem this thesis targets on GPFS is not an
artefact of one filesystem.

The many-small-file problem is a specific and well-known pathology on parallel
filesystems.
Metadata cost per file is close to constant regardless of how few bytes are read,
so per-file overhead does not amortise, and a corpus of thousands of small images
generates metadata traffic out of proportion to the data actually transferred.
The standard mitigation is to aggregate samples into large sequentially read
archives.
WebDataset~\cite{Aizman19} formalises this: samples are stored in POSIX tar archives
and streamed sequentially, converting many small random reads into few large
sequential ones, at the cost of exact random-access shuffling - which is
approximated by shuffling shard order and maintaining an in-memory sample buffer.
NVIDIA DALI~\cite{DALI} attacks an adjacent part of the problem by moving decode and
augmentation onto the GPU.
The term \emph{IO-aware}, used throughout this thesis for \pads{}'s central
mechanism, echoes an argument made one level lower in the compute stack: Dao et
al.~\cite{Dao22} show that attention kernels which account explicitly for the cost
of moving data between GPU high-bandwidth memory and on-chip SRAM outperform
kernels that minimise floating-point operations alone.
The claim in this thesis is the same principle applied one layer higher in the
storage hierarchy - between a shared parallel filesystem and node-local scratch,
rather than between HBM and SRAM.

Two gaps in this body of work motivate \pads{}.

First, these systems are \emph{statically configured}.
WebDataset streams shards at a fixed depth; the number of loader workers and the
prefetch factor are chosen once.
None of them observes whether the configuration is in fact sufficient to hide the
transfer, and none revises it when contention on a shared filesystem changes the
answer mid-run.

Second, and more specifically, they do not exploit the fact that in supervised
training with a fixed epoch permutation, \emph{future accesses are known exactly}.
Caching and prefetching are ordinarily speculative because future references must be
predicted.
Here they need not be: once the permutation is drawn, the identity of the next $k$
samples is determined, and a staging schedule can be computed rather than guessed.
\pads{} is built on that observation.

\subsection{Summary}

Segmentation architecture for this workload is a solved problem in the sense that
matters here: the reference study found that architectural choice moves IoU by only
a few percentage points, plausibly within the fluctuation induced by random
augmentation~\cite{Casini23}.
If the remaining headroom is not in the model, it may be in the pipeline that feeds
it.
The I/O literature establishes that data stalls are real and that aggregation into
shards mitigates the dominant filesystem pathology, but leaves the configuration
static and the prefetch speculative.
This thesis addresses that gap.
